\section{Praktikum}
\label{sec:praktikum-bab15}

\begin{praktikumbox}
\textbf{Judul:} \eng{Debugging} terpandu tiga kasus bug tingkat
instruksi.\\[0.3em]
\textbf{Sub-CPMK:} 4.1 dan 4.2 \quad\textbullet\quad
\textbf{Alokasi waktu:} 100 menit \quad\textbullet\quad
\textbf{Bentuk:} individu, laporan tertulis
\end{praktikumbox}

\subsection{Tujuan Praktikum}

Setelah menyelesaikan praktikum ini, mahasiswa mampu menelusuri
program pada level instruksi dengan perintah \code{R}, \code{U},
\code{T}, \code{P}, \code{G}, dan \code{D}; mengisi tabel \eng{trace}
register dan \eng{flag}; membaca \eng{dump} memori dan \eng{stack};
serta menemukan dan memperbaiki perulangan tak berujung, kesalahan
\eng{off-by-one}, dan \eng{stack} yang tidak seimbang.

\subsection{Alat dan Bahan}

\begin{itemize}
    \item Komputer dengan sistem operasi Windows.
    \item DOSBox beserta berkas \file{DEBUG.EXE}.
    \item TASM dan TLINK (atau Emu8086 sebagai pembanding).
    \item Berkas \file{bug-loop.asm}, \file{bug-stack.asm}, dan
          \file{program-larik.asm} yang disiapkan dosen.
\end{itemize}

\subsection{Langkah Kerja}

\begin{enumerate}
    \item \textbf{Menyiapkan berkas uji.}
    \begin{enumerate}
        \item Salin berkas \file{bug-loop.asm} ke direktori kerja
              DOSBox, lalu rakit dan gabungkan dengan perintah
              \file{tasm bug-loop.asm} dan \file{tlink bug-loop.obj}.
        \item Ulangi langkah yang sama untuk \file{bug-stack.asm}.
        \item Tuliskan program \eng{off-by-one} pada
              Kode~\ref{lst:off-by-one} ke dalam berkas
              \file{program-larik.asm}, lalu rakit.
    \end{enumerate}

    \item \textbf{Kasus 1: perulangan tak berujung
          (\file{bug-loop.exe}).}
    \begin{enumerate}
        \item Jalankan DEBUG, lalu muat program dengan urutan
              \code{N} dan \code{L}.
        \item Catat keluaran \code{R}. Berapa nilai \reg{CS}, \reg{IP},
              \reg{DS}, \reg{SS}, dan \reg{SP}?
        \item Tampilkan hasil pembongkaran dengan \code{U 0000 001c},
              lalu salin ke dalam laporan.
        \item Jalankan \code{T} sekurang-kurangnya sepuluh kali.
              Catat setiap langkah pada tabel \eng{trace}
              (gunakan susunan kolom pada
              Tabel~\ref{tab:trace-register}).
        \item Jawab: pada langkah berapa penyimpangan mulai terjadi,
              dan register apa yang menyimpang?
        \item Tuliskan hipotesis penyebab, buktikan dengan data
              \eng{trace}, lalu ubah program dan rakit ulang sampai
              keluarannya benar (\hex{000F} atau \dec{15}).
    \end{enumerate}

    \item \textbf{Kasus 2: kesalahan \eng{off-by-one}
          (\file{program-larik.exe}).}
    \begin{enumerate}
        \item Muat program ke dalam DEBUG, lalu jalankan sampai
              perulangan selesai menggunakan \code{G} pada alamat
              setelah \asm{LOOP}.
        \item Tampilkan isi larik dengan \code{D DS:0}.
        \item Tentukan berapa byte yang seharusnya dibaca dan berapa
              byte yang benar-benar dibaca.
        \item Perbaiki program, rakit ulang, dan buktikan bahwa
              hasilnya sekarang \hex{0096}.
    \end{enumerate}

    \item \textbf{Kasus 3: \eng{stack} tidak seimbang
          (\file{bug-stack.exe}).}
    \begin{enumerate}
        \item Muat program, lalu jalankan dengan \code{G 0005} agar
              berhenti tepat sebelum \asm{CALL}.
        \item Telusuri dengan \code{T} sampai seluruh register selesai
              ditekan. Perhatikan nilai \reg{SP} pada setiap langkah
              dan catat pada tabel \eng{trace} dengan menambahkan
              kolom \reg{SP}.
        \item Tampilkan \eng{stack} dengan \code{D SS:SP}, lalu
              gambarkan susunan bingkai \eng{stack}-nya seperti pada
              Gambar~\ref{fig:stack-frame}.
        \item Lanjutkan penelusuran sampai \asm{RET} dijalankan.
              Alamat berapa yang diambil oleh \asm{RET}, dan mengapa
              demikian?
        \item Bandingkan pemakaian \code{T} dan \code{P} pada
              instruksi \asm{CALL} tersebut. Tuliskan perbedaan
              keluarannya dan jelaskan kapan masing-masing lebih
              tepat dipakai.
        \item Perbaiki program sehingga pesan \code{Selesai.}
              tercetak.
    \end{enumerate}
\end{enumerate}

\subsection{Tabel Kerja}

Salin dan lengkapi tabel berikut pada jurnal praktikum Anda.

\begin{table}[H]
    \centering
    \caption{Templat tabel \eng{trace} untuk diisi mahasiswa.}
    \label{tab:template-trace}
    \scriptsize
    \setlength{\tabcolsep}{2.5pt}
    \begin{tabular}{@{}rl l rrrr r ccc p{2.3cm}@{}}
        \toprule
        \textbf{L} & \textbf{Alamat} & \textbf{Instruksi} &
        \textbf{AX} & \textbf{BX} & \textbf{CX} & \textbf{DX} &
        \textbf{SP} & \textbf{ZF} & \textbf{CF} & \textbf{OF} &
        \textbf{Keterangan} \\
        \midrule
        1 & & & & & & & & & & & \\
        \addlinespace
        2 & & & & & & & & & & & \\
        \addlinespace
        3 & & & & & & & & & & & \\
        \addlinespace
        4 & & & & & & & & & & & \\
        \addlinespace
        5 & & & & & & & & & & & \\
        \bottomrule
    \end{tabular}
\end{table}

\subsection{Hasil yang Dikumpulkan}

\begin{enumerate}
    \item Tabel \eng{trace} lengkap untuk ketiga kasus, minimal
          sepuluh langkah per kasus, memuat kolom \reg{AX},
          \reg{BX}, \reg{CX}, \reg{DX}, \reg{SP}, \flag{ZF},
          \flag{CF}, dan \flag{OF}.
    \item Salinan keluaran \code{R}, \code{U}, \code{D}, dan
          \code{T} untuk setiap kasus.
    \item Rumusan hipotesis dan bukti pengujiannya untuk setiap
          kasus, ditulis dengan alur metodologi pada
          Gambar~\ref{fig:metodologi-debugging}.
    \item Berkas hasil perbaikan ketiga program beserta tangkapan
          layar keluarannya.
    \item Perbandingan pengalaman memakai \code{T} dan \code{P} pada
          instruksi \asm{CALL}.
\end{enumerate}

\begin{asesmenbox}
\textbf{Asesmen kinerja \eng{debugging} (bobot 10\% dari komponen
praktikum).} Penilaian dilakukan terhadap laporan dan demonstrasi
langsung.

\begin{tabular}{@{}p{2.6cm}cp{9.6cm}@{}}
    \toprule
    \textbf{Aspek kinerja} & \textbf{Bobot} & \textbf{Kriteria} \\
    \midrule
    Penggunaan perintah DEBUG & 20\% &
    Mampu memuat berkas (\code{N}, \code{L}), menampilkan register
    (\code{R}), membongkar kode (\code{U}), dan menelusuri
    (\code{T}, \code{P}, \code{G}) tanpa bantuan \\
    \addlinespace
    Ketepatan tabel \eng{trace} & 25\% &
    Nilai register dan \eng{flag} tercatat benar pada setiap langkah;
    perbedaan \code{T} dan \code{P} dicatat dengan tepat \\
    \addlinespace
    Analisis \eng{dump} memori & 20\% &
    Mampu menafsirkan isi larik dan bingkai \eng{stack}; menemukan
    byte atau alamat yang menyimpang \\
    \addlinespace
    Ketepatan diagnosis & 20\% &
    Hipotesis sesuai bukti; penyebab yang disebutkan benar-benar
    sumber kesalahan, bukan gejala \\
    \addlinespace
    Mutu perbaikan dan verifikasi & 15\% &
    Perbaikan mengubah satu hal pada satu waktu dan seluruh kasus uji
    diulang setelah perbaikan \\
    \bottomrule
\end{tabular}
\end{asesmenbox}
